\scanpage{37}
\mspage{28}
% archive sheet number: 37
% inv: deletion | top | Il semble qu'il manque
\deleted{Il \unsure{semble} qu'il manque une condition pour pouvoir le conclure !}

% inv: text | top | Il semble qu'on ait besoin
Il semble qu'on ait besoin aussi de la condition
% inv: text | upper third | d) forall X in M
 d) $\forall X\in M$, $\varphi$ induit une \uline{bijection}
% inv: display | upper third | tilde X to tilde X prime
\[
\widetilde X\xrightarrow{\sim}\widetilde X'\qquad\text{où}\quad X'=\varphi(X)
\]
% inv: margin note | slanted left upper | dated condition d)
% margin note written sideways
% slanted
\margin{\unsure{12|11}.7.86\par Cette condition d) devrait \unsure{rester|être} \unsure{facultative}. Il y a des ``\unsure{sous-mag.}'' qui ne sont pas fermés par $\unsure{\leq|\Delta}$ !}

% inv: text | upper middle | Corollaire 1
\uline{Corollaire 1} Les conditions a), c), d) impliquent que $\forall F\in\mathrm{Fig}(M)$, on a $\varphi(F)\in\mathrm{Fig}(M')$, et $\varphi|F$ est injective \inserted{et $\simeq$ d'ens. ordonnés}.

% inv: text | middle | Immédiat par la proposition
Immédiat par la proposition précédente \inserted{(si on pose $I=\widetilde F$)} si $\varphi$ est injective. \deleted{Dans le cas général,} Il suffit de vérifier que $\varphi|F$ est injective. Mais si $X,Y\in F$, \deleted{\unsure{si}} $X\ne Y$, on aura $X\mathrel{|\circ|}Y$, donc $\varphi(X)\mathrel{|\circ|}\varphi(Y)$ par c), donc $\varphi(X)\ne\varphi(Y)$ OK ! \deleted{\unsure{Montrons que} $\varphi$} \inserted{C'est la prop qui} montre de plus que $\varphi$ est un isom. ordonné\ldots

% inv: text | middle | Corollaire 2
\uline{Corollaire 2}\quad\deleted{Si $X\unsure{\gtrless|\bowtie}Y$} Si $X'=\varphi(X)$, $Y'=\varphi(\unsure{Y})$
% inv: display | middle | compatibility and parallel preserved
\[
\begin{aligned}
X\gtrless Y&\Longrightarrow X'\gtrless Y'\\
X\parallel Y&\Longrightarrow X'\parallel Y'.
\end{aligned}
\]
% inv: text | middle lower | En effet X relation Y
En effet, $X\gtrless Y$ signifie que $X,Y$ font partie d'une figure \inserted{$F$}, et on applique Cor. 1. Plus \deleted{Si $X\parallel Y$, \unsure{on a} $X'$ \unsure{on aura} $\mathrm{Inf}(X,Y)$} précisément

% inv: text | lower third | Cor 3
\uline{Cor 3}\quad Soit $F$ figure dans $M$. Alors pour $X,Y\in\widetilde F$, et $X'=\varphi(X)$, $Y'=\varphi(Y)$, on a
% inv: display | lower third | deleted relations then two equivalences
\[
\begin{aligned}
&\deleted{X\gtrless Y\Longleftrightarrow X'\gtrless Y'}\\
&\deleted{X\parallel Y}\\
&X\leq Y\Longleftrightarrow X'\leq Y'\\
&\deleted{\unsure{X\gtrless Y}}\\
&X\parallel Y\Longleftrightarrow X'\parallel Y'
\end{aligned}
\]
% inv: margin note | right lower | se généralise
\margin{\unsure{se généralise} aux sous-figures quelconques de $F$}

% inv: text | bottom | La première
La première\inserted{$\unsure{*}$} \unsure{équivalence} \deleted{\unsure{provient}} déjà \unsure{vue|vu} (Cor. 1). La deuxième provient du fait que $X\parallel Y$ se lit dans l'ens ordonné $I=\widetilde F$ ($\{X,Y\}$ non minoré) lequel est isom. à $I'=\varphi(I)$.
